\documentclass[UTF8,12pt,a4paper,fontset=fandol]{ctexart}

% 支持从项目根目录或当前笔记目录编译。
\makeatletter
\def\input@path{{../}{../../}}
\makeatother
\usepackage{researchnotes}

\title{矩阵常见运算笔记}
\author{}
\date{}

\begin{document}
\maketitle

\section{从矩阵的尺寸开始}

矩阵既能组织数据，也能表示线性变换。矩阵运算沿用了加、减、乘等符号，
但运算能否进行、结果的尺寸以及因子的顺序，都需要单独检查。
本笔记从加减和乘法出发，介绍转置、幂、行列式与求逆，再讨论消元、秩、分块运算及范数。
这些运算的输出可能是矩阵，也可能是一个数；理解输出的类型，是正确使用公式的第一步。

除特别说明外，全文讨论实矩阵。记 $A=(a_{ij})\in\mathbb R^{m\times n}$，表示 $A$ 有 $m$ 行、$n$ 列，
$a_{ij}$ 是第 $i$ 行、第 $j$ 列的元素，其中 $1\le i\le m$、$1\le j\le n$。
各尺寸均为正整数，向量默认写成列向量；复数情形只在介绍共轭转置时另行说明。
行数与列数相等的矩阵称为\keyterm{方阵}（square matrix）。两个矩阵相等，
是指它们尺寸相同且对应元素全部相等。
\keyterm{零矩阵}（zero matrix）$0_{m\times n}$ 的元素全为零；
\keyterm{单位矩阵}（identity matrix）$I_n$ 的主对角线元素为 $1$，其余为 $0$。

\section{加减法、数乘与除以非零数}

\begin{definition}[逐项线性运算]
设 $A=(a_{ij})$、$B=(b_{ij})$ 都属于 $\mathbb R^{m\times n}$，$\lambda\in\mathbb R$。
\keyterm{矩阵加法}（matrix addition）、\keyterm{矩阵减法}（matrix subtraction）
与\keyterm{数乘}（scalar multiplication）分别定义为
\[
  A+B=(a_{ij}+b_{ij}),\qquad A-B=(a_{ij}-b_{ij}),\qquad
  \lambda A=(\lambda a_{ij}).
\]
若 $\lambda\ne0$，定义 $A/\lambda=(1/\lambda)A$，即每个元素都除以 $\lambda$。
\end{definition}

加减法要求两个矩阵\keyterm{尺寸完全相同}，结果尺寸不变。
由于这些运算逐项进行，实数运算律给出 $A+B=B+A$、$(A+B)+C=A+(B+C)$，
以及 $\lambda(A+B)=\lambda A+\lambda B$；其中 $C$ 与 $A,B$ 同型。
更一般地，同型矩阵 $A_1,\ldots,A_r$ 与实数 $\lambda_1,\ldots,\lambda_r$ 可以组成
\keyterm{线性组合}（linear combination）$\lambda_1A_1+\cdots+\lambda_rA_r$。
普通矩阵代数中，矩阵与实数直接相加没有默认定义；对方阵写 $A+\lambda I_n$，
表示只给主对角线元素加上 $\lambda$，并不是给所有元素加上 $\lambda$。

\begin{example}[贯穿全文的一组矩阵]
\label{ex:matrix-operations-pair}
取 $A=\begin{pmatrix}1&2\\0&1\end{pmatrix}$、
$B=\begin{pmatrix}2&0\\1&3\end{pmatrix}$。逐项计算得到
\[
  A+B=\begin{pmatrix}3&2\\1&4\end{pmatrix},\qquad
  A-B=\begin{pmatrix}-1&2\\-1&-2\end{pmatrix},\qquad
  2A=\begin{pmatrix}2&4\\0&2\end{pmatrix}.
\]
例如，$(A-B)_{12}=2-0=2$。数乘中的 $2$ 是一个实数，并非二阶单位矩阵。
\end{example}

\section{矩阵乘法：行与列配对求和}

\subsection{定义、尺寸与运算律}

\begin{definition}[矩阵乘法]
若 $A\in\mathbb R^{m\times n}$、$B\in\mathbb R^{n\times p}$，则
\keyterm{矩阵乘积}（matrix product）$C=AB\in\mathbb R^{m\times p}$ 定义为
\begin{equation}
  c_{ij}=\sum_{k=1}^{n}a_{ik}b_{kj},\qquad
  1\le i\le m,\quad 1\le j\le p.
  \label{eq:matrix-product-entry}
\end{equation}
\end{definition}

即取 $A$ 的第 $i$ 行与 $B$ 的第 $j$ 列，对应相乘后求和。
尺寸规则是 $(m\times n)(n\times p)\longrightarrow m\times p$：
\keyterm{左矩阵的列数必须等于右矩阵的行数}。
若 $x\in\mathbb R^p$ 是列向量，则 $(AB)x=A(Bx)$，所以 $AB$ 表示先施加 $B$、再施加 $A$。

对例题~\ref{ex:matrix-operations-pair} 中的矩阵，按式~\eqref{eq:matrix-product-entry} 得
\[
  AB=\begin{pmatrix}1\cdot2+2\cdot1&1\cdot0+2\cdot3\\
                     0\cdot2+1\cdot1&0\cdot0+1\cdot3\end{pmatrix}
     =\begin{pmatrix}4&6\\1&3\end{pmatrix},\qquad
  BA=\begin{pmatrix}2&4\\1&5\end{pmatrix}.
\]
这里 $AB\ne BA$，说明矩阵乘法通常\keyterm{不满足交换律}。
对一般长方形矩阵，$AB$ 有定义甚至不保证 $BA$ 有定义。
在各项尺寸相容时，仍有结合律 $(AB)C=A(BC)$、分配律
$A(B+C)=AB+AC$、$(A+B)C=AC+BC$，且 $I_mA=A=AI_n$。
例如，分配律可由 $\sum_k a_{ik}(b_{kj}+c_{kj})
=\sum_k a_{ik}b_{kj}+\sum_k a_{ik}c_{kj}$ 逐项核查。

\subsection{哈达玛积：逐元素相乘}
\label{subsec:matrix-hadamard-product}

\begin{definition}[哈达玛积]
设 $A=(a_{ij})$、$B=(b_{ij})$ 都属于 $\mathbb R^{m\times n}$。
二者的\keyterm{哈达玛积}（Hadamard product），也称\keyterm{舒尔积}（Schur product）
或\keyterm{逐元素乘积}（element-wise product），定义为
\begin{equation}
  A\odot B=(a_{ij}b_{ij})\in\mathbb R^{m\times n},\qquad
  (A\odot B)_{ij}=a_{ij}b_{ij}.
  \label{eq:matrix-hadamard-product}
\end{equation}
\end{definition}

它只把相同位置的两个元素相乘，\keyterm{不对行、列配对求和}。
两个因子必须同型，但不必是方阵，结果尺寸也不变。
例如，两个 $2\times3$ 矩阵可以计算哈达玛积，却不能直接计算普通矩阵乘积。
本文用 $A\odot B$ 表示哈达玛积，用 $AB$ 表示普通矩阵乘积；两种记号不能混用。

\begin{example}[与普通矩阵乘积比较]
对例题~\ref{ex:matrix-operations-pair} 中的 $A,B$，逐项相乘得
\[
  A\odot B
  =\begin{pmatrix}1\cdot2&2\cdot0\\0\cdot1&1\cdot3\end{pmatrix}
  =\begin{pmatrix}2&0\\0&3\end{pmatrix},\qquad
  AB=\begin{pmatrix}4&6\\1&3\end{pmatrix}.
\]
例如，$(A\odot B)_{12}=a_{12}b_{12}=0$，而
$(AB)_{12}=a_{11}b_{12}+a_{12}b_{22}=6$。
前者只依赖一个位置，后者汇总一行与一列的配对乘积。
\end{example}

\begin{property}[哈达玛积的运算律]
对任意同型实矩阵 $A,B,C$ 和实数 $\lambda$，有
\begin{align*}
  A\odot B&=B\odot A,&
  (A\odot B)\odot C&=A\odot(B\odot C),\\
  A\odot(B+C)&=A\odot B+A\odot C,&
  (\lambda A)\odot B&=\lambda(A\odot B).
\end{align*}
\end{property}

这些性质都由实数逐项运算得到。例如，分配律两边的第 $(i,j)$ 个元素分别是
$a_{ij}(b_{ij}+c_{ij})$ 与 $a_{ij}b_{ij}+a_{ij}c_{ij}$，故相等。
特别地，哈达玛积总满足交换律，这与普通矩阵乘法不同。

记 $J_{m\times n}$ 为所有元素均为 $1$ 的矩阵，则 $A\odot J_{m\times n}=A$，
所以哈达玛积的\keyterm{单位元}（identity element）是全 $1$ 矩阵。
对于 $n$ 阶方阵，$A\odot I_n$ 只保留 $A$ 的主对角线元素，其余位置变为零，通常不等于 $A$。
更一般地，若 $M=(m_{ij})$ 的元素均为 $0$ 或 $1$，则 $A\odot M$ 在 $m_{ij}=1$ 的位置保留原值，
在 $m_{ij}=0$ 的位置置零；这种 $M$ 称为\keyterm{掩码矩阵}（mask matrix）。
若各位置需要不同的权重 $w_{ij}$，则用同型权重矩阵 $W=(w_{ij})$ 计算 $W\odot A$。

与之配套的\keyterm{逐元素除法}（element-wise division）定义为
$A\oslash B=(a_{ij}/b_{ij})$，要求每个 $b_{ij}\ne0$。
令 $B^{\odot(-1)}=(1/b_{ij})$，则
\[
  B\odot B^{\odot(-1)}=J_{m\times n},\qquad
  A\oslash B=A\odot B^{\odot(-1)}.
\]
因此，对给定的同型矩阵 $B,C$，若 $B$ 的每个元素都非零，方程 $B\odot X=C$ 的唯一解是
$X=C\oslash B$。若某个 $b_{ij}=0$，则相应位置要求 $c_{ij}=0$ 才可能有解，
且该位置的 $x_{ij}$ 可以任取。
这里的逐元素倒数 $B^{\odot(-1)}$ 与普通逆矩阵 $B^{-1}$ 是不同概念，后者将在求逆一节定义。

\subsection{矩阵与向量、内积与外积}

若将 $A\in\mathbb R^{m\times n}$ 的列写成 $a_1,\ldots,a_n$，则对
$x=(x_1,\ldots,x_n)^{\mathsf T}$，有 $Ax=x_1a_1+\cdots+x_na_n$。
因此矩阵乘向量就是按向量分量对矩阵各列作线性组合。
若 $B$ 的各列为 $b_1,\ldots,b_p$，便有 $AB=(Ab_1,\ldots,Ab_p)$。

对 $u,v\in\mathbb R^n$，\keyterm{内积}（inner product）$u^{\mathsf T}v=\sum_i u_iv_i$
是一个实数；对 $u\in\mathbb R^m$、$v\in\mathbb R^n$，
\keyterm{外积}（outer product）$uv^{\mathsf T}=(u_iv_j)$ 是 $m\times n$ 矩阵。
这里 $\mathsf T$ 表示下一节介绍的转置，它将列向量变为行向量。例如
\[
  u=\begin{pmatrix}1\\2\end{pmatrix},\quad
  v=\begin{pmatrix}3\\4\end{pmatrix},\qquad
  u^{\mathsf T}v=11,\qquad uv^{\mathsf T}=\begin{pmatrix}3&4\\6&8\end{pmatrix}.
\]

矩阵乘法也不能无条件约去因子。取
$P=\begin{pmatrix}1&0\\0&0\end{pmatrix}$、$Q=\begin{pmatrix}0&0\\0&1\end{pmatrix}$，
则 $P,Q$ 均非零，却有 $PQ=0_{2\times2}$。
因此，$AB=AC$ 一般不能推出 $B=C$；若 $A$ 可逆，则可以左乘 $A^{-1}$ 约去 $A$。

\section{转置与矩阵的幂}

\subsection{转置与共轭转置}

\keyterm{转置}（transpose）把行、列互换：若 $A\in\mathbb R^{m\times n}$，
则 $A^{\mathsf T}\in\mathbb R^{n\times m}$，且 $(A^{\mathsf T})_{ij}=a_{ji}$。
在运算有定义时，有
\[
  (A^{\mathsf T})^{\mathsf T}=A,\qquad
  (A+B)^{\mathsf T}=A^{\mathsf T}+B^{\mathsf T},\qquad
  (AB)^{\mathsf T}=B^{\mathsf T}A^{\mathsf T}.
\]
最后一式会颠倒因子顺序，因为两边第 $(i,j)$ 个元素分别为
$\sum_k a_{jk}b_{ki}$ 与 $\sum_k b_{ki}a_{jk}$，二者相等。
数乘也满足 $(\lambda A)^{\mathsf T}=\lambda A^{\mathsf T}$。
满足 $A^{\mathsf T}=A$ 的实方阵称为\keyterm{对称矩阵}（symmetric matrix）；
对任意实矩阵 $A$，$A^{\mathsf T}A$ 都是对称方阵。

若 $A$ 的元素为复数，则常用\keyterm{共轭转置}（conjugate transpose）
$A^*=\overline{A}^{\mathsf T}$，即先逐项取复共轭，再转置。
例如 $z=(1,\mathrm i)^{\mathsf T}$（$\mathrm i^2=-1$）满足
$z^{\mathsf T}z=0$，而 $z^*z=2$。
复向量的标准内积使用 $u^*v$，以保证 $z^*z=\sum_i|z_i|^2$ 非负且仅在 $z=0$ 时为零。
对实矩阵，共轭转置与转置相同。

\subsection{整数次幂与多项式运算}

对于 $n$ 阶方阵 $A$，\keyterm{矩阵幂}（matrix power）定义为
$A^0=I_n$、$A^r=\underbrace{A\cdots A}_{r\text{ 个因子}}$（$r$ 为正整数）。
结合律给出 $A^rA^s=A^{r+s}$，其中 $r,s$ 为非负整数。
继续使用例题~\ref{ex:matrix-operations-pair} 中的 $A$，有
\[
  A^{\mathsf T}=\begin{pmatrix}1&0\\2&1\end{pmatrix},\qquad
  A^2=\begin{pmatrix}1&4\\0&1\end{pmatrix},\qquad
  A^3=\begin{pmatrix}1&6\\0&1\end{pmatrix}.
\]
矩阵幂不等于逐元素取幂：这里 $(A^3)_{12}=6$，而 $a_{12}^3=2^3=8$。
它也不同于数乘，例如 $A^2$ 的主对角线元素为 $1$，而 $2A$ 的主对角线元素为 $2$。

对同阶方阵，分配律给出
\[
  (A+B)^2=A^2+AB+BA+B^2.
\]
只有在 $AB=BA$ 时，才能将两个中间项合并为 $2AB$。
同样，$(AB)^r=A^rB^r$ 对一般矩阵不成立；当 $AB=BA$ 时，
可通过交换相邻的 $A,B$ 并利用结合律，对所有非负整数 $r$ 得到该式。
若 $p(t)=c_0+c_1t+\cdots+c_rt^r$，则定义
\keyterm{矩阵多项式}（matrix polynomial）$p(A)=c_0I_n+c_1A+\cdots+c_rA^r$，
其中常数项必须写成 $c_0I_n$。

\section{行列式与迹：从方阵得到一个数}
\label{sec:matrix-determinant-trace}

\subsection{行列式的计算}

\keyterm{行列式}（determinant）是与方阵对应的一个数，记为 $\det A$。
可沿第一行递归定义：一阶情形 $\det(a)=a$；当 $n\ge2$ 时，
令 $M_{1j}$ 为删去 $A$ 的第一行和第 $j$ 列所得的子矩阵，则
\[
  \det A=\sum_{j=1}^n(-1)^{1+j}a_{1j}\det M_{1j},\qquad
  \det\begin{pmatrix}a&b\\c&d\end{pmatrix}=ad-bc.
\]
这里的交错符号不能省略，行列式也不是把全部元素相乘。
三角矩阵的行列式等于主对角线元素之积；例如上三角矩阵可沿第一列递归展开得到这一结论。
因此，例题~\ref{ex:matrix-operations-pair} 中 $\det A=1$，$\det B=6$。

对同阶方阵，下列是行列式的基本运算性质：
\[
  \det(AB)=\det A\det B,\qquad
  \det(A^{\mathsf T})=\det A,\qquad
  \det(\lambda A)=\lambda^n\det A.
\]
最后一式中的 $n$ 是矩阵阶数，因为每一行都被乘以 $\lambda$。
行列式对各行分别线性，但不是整个矩阵的线性函数：一般
$\det(A+B)\ne\det A+\det B$，例如二阶情形 $\det(2I_2)=4\ne2$。
本节列出计算时使用的基本性质，不展开一般阶行列式理论的完整证明。

\subsection{迹与循环换位}

方阵的\keyterm{迹}（trace）是主对角线元素之和：
$\operatorname{tr}A=\sum_{i=1}^n a_{ii}$。
由求和定义，$\operatorname{tr}(\lambda A+\mu B)
=\lambda\operatorname{tr}A+\mu\operatorname{tr}B$，其中 $A,B$ 同阶、$\lambda,\mu\in\mathbb R$。
对 $A\in\mathbb R^{m\times n}$、$B\in\mathbb R^{n\times m}$，还有
\[
  \operatorname{tr}(AB)=\sum_{i=1}^m\sum_{j=1}^n a_{ij}b_{ji}
  =\operatorname{tr}(BA).
\]
即使 $AB\ne BA$，二者的迹仍相等。对前述二阶例子，两者的迹都是 $7$。
对多个因子，在乘积有定义且为方阵时可作循环换位，如
$\operatorname{tr}(ABC)=\operatorname{tr}(BCA)$，但不能据此任意交换两个因子。

\section{逆矩阵与“矩阵除法”}

\subsection{求逆的条件与公式}

\begin{definition}[逆矩阵]
对 $n$ 阶方阵 $A$，若存在 $n$ 阶矩阵 $D$ 使 $AD=DA=I_n$，则称 $A$
\keyterm{可逆}（invertible），称 $D$ 为其\keyterm{逆矩阵}（inverse matrix），记为 $A^{-1}$。
\end{definition}

逆矩阵若存在便唯一：若 $D,E$ 都是逆矩阵，则 $D=D(AE)=(DA)E=E$。
并非每个方阵都可逆，例如零矩阵与任意矩阵相乘仍为零，不可能得到单位矩阵。
求逆也不是对元素取倒数。对前面的 $A$，直接乘法验证可得
\[
  A^{-1}=\begin{pmatrix}1&-2\\0&1\end{pmatrix},\qquad
  AA^{-1}=A^{-1}A=I_2.
\]

方阵可逆当且仅当其行列式非零。其计算依据是第~\ref{sec:matrix-elimination-rank} 节介绍的消元：
初等行变换不改变可逆性及行列式是否为零，
而阶梯形方阵可逆恰好要求每列都有主元，此时对角线元素之积非零。
二阶矩阵尤其有
\begin{equation}
  \begin{pmatrix}a&b\\c&d\end{pmatrix}^{-1}
  =\frac{1}{ad-bc}\begin{pmatrix}d&-b\\-c&a\end{pmatrix},
  \qquad ad-bc\ne0.
  \label{eq:matrix-inverse-two-by-two}
\end{equation}
原矩阵与右侧分子矩阵的两个顺序的乘积均为 $(ad-bc)I_2$，故公式可直接验证。
可逆矩阵满足 $(A^{\mathsf T})^{-1}=(A^{-1})^{\mathsf T}$，
由 $AA^{-1}=A^{-1}A=I_n$ 两边转置即可核查。

\subsection{通过矩阵方程理解除法}

由于乘法的顺序不能任意交换，通常不使用含义不明的分式 $B/A$。
当 $A$ 是可逆的 $n$ 阶方阵时，应通过方程说明要求什么：
\begin{equation}
  AX=B\ \Longleftrightarrow\ X=A^{-1}B,\qquad
  YA=C\ \Longleftrightarrow\ Y=CA^{-1},
  \label{eq:matrix-equations-inverse}
\end{equation}
其中 $B,X\in\mathbb R^{n\times p}$，$C,Y\in\mathbb R^{q\times n}$。
第一式在两边左乘 $A^{-1}$，第二式在两边右乘 $A^{-1}$；代回即验证存在性与唯一性。
例如对前述 $A,B$，方程 $AX=B$ 的唯一解为
\[
  X=A^{-1}B=\begin{pmatrix}0&-6\\1&3\end{pmatrix},\qquad
  AX=\begin{pmatrix}2&0\\1&3\end{pmatrix}=B.
\]

可逆矩阵还可定义负整数次幂 $A^{-r}=(A^{-1})^r$（$r$ 为正整数）。
若 $A,B$ 都是同阶可逆方阵，则 $(AB)^{-1}=B^{-1}A^{-1}$：
按此顺序相乘，两侧的乘积都能依次消去相邻的互逆因子，得到单位矩阵。

第~\ref{subsec:matrix-hadamard-product} 小节定义的逐元素除法 $A\oslash B$，
不等于这里通过普通矩阵乘法求逆或解方程。例如全为 $1$ 的二阶矩阵可作逐元素除法的除数，
但它不可逆：它将非零列向量 $(1,-1)^{\mathsf T}$ 映为零，若有逆矩阵便会推出该向量为零，产生矛盾。

\section{初等行变换、秩与线性方程组}
\label{sec:matrix-elimination-rank}

\keyterm{初等行变换}（elementary row operation）包括三种操作：交换两行；
将一行乘以非零数；将一行的某个倍数加到另一行。
它们都有逆操作，所以不会改变各行所含的线性约束信息。
求解 $Ax=b$ 时，必须同时变换系数与右端，即对\keyterm{增广矩阵}
（augmented matrix）$[A\mid b]$ 进行行变换。
行变换对行列式的影响需要另记：交换两行使行列式变号；一行乘以 $c\ne0$，行列式也乘以 $c$；
将一行的倍数加到另一行，行列式不变。因此可通过消元化为三角矩阵，并追踪这些因子来计算行列式。

\keyterm{高斯消元}（Gaussian elimination）逐列寻找非零元素作为
\keyterm{主元}（pivot），必要时交换行，再消去其下方元素；某列没有可用非零元素时跳过该列。
处理完各列后得到\keyterm{行阶梯形}（row echelon form）：零行位于底部，
每个非零行的首个非零元素位于上一行首个非零元素的右侧。
继续将主元归一化并消去主元上方元素，可得\keyterm{简化行阶梯形}
（reduced row echelon form）。

矩阵的\keyterm{秩}（rank）$\operatorname{rank}A$ 定义为其列向量组中线性无关向量的最大个数。
行变换相当于左乘可逆的初等矩阵，保持列向量之间的线性关系，因而保持秩。
阶梯形矩阵的主元列线性无关，其余列可由主元列线性表示，
所以秩等于主元个数，也等于阶梯形矩阵的非零行数。
特别地，$0\le\operatorname{rank}A\le\min(m,n)$。

\begin{example}[用消元同时求秩与解方程]
求 $R=\begin{pmatrix}1&2&1\\2&4&3\end{pmatrix}$ 的秩，并解
$Rx=(4,9)^{\mathsf T}$。
\end{example}
\begin{solve}
用第 $1$ 行的 $-2$ 倍加到第 $2$ 行，再用第 $2$ 行的 $-1$ 倍加到第 $1$ 行，得到
\[
  \left(\begin{array}{ccc|c}1&2&1&4\\2&4&3&9\end{array}\right)
  \longrightarrow
  \left(\begin{array}{ccc|c}1&2&1&4\\0&0&1&1\end{array}\right)
  \longrightarrow
  \left(\begin{array}{ccc|c}1&2&0&3\\0&0&1&1\end{array}\right).
\]
系数部分有两个主元，故 $\operatorname{rank}R=2$。
第二个未知量是自由变量，令 $x_2=t$，则全部解为
$x=(3-2t,t,1)^{\mathsf T}$，$t\in\mathbb R$。
该矩阵有三列而只有两个主元，因此本例有无穷多解。
\end{solve}

一般地，$Ax=b$ 有解当且仅当 $\operatorname{rank}A=\operatorname{rank}[A\mid b]$：
秩不等时，消元会出现 $0=\beta$（$\beta\ne0$）的矛盾行。
在有解的前提下，若 $A$ 有 $n$ 列且秩为 $n$，则解唯一；若秩小于 $n$，则存在自由变量，实数范围内有无穷多解。

求逆也可用同样的操作：将 $[A\mid I_n]$ 化为 $[I_n\mid A^{-1}]$。
若所有行变换的乘积为 $E$，则左右两块分别变为 $EA$ 与 $E$；当 $EA=I_n$ 时，$E=A^{-1}$。
若左块无法化为 $I_n$，则 $A$ 不可逆。
实际求解 $AX=B$ 时，可以直接对 $[A\mid B]$ 消元，从而同时求出 $X$ 的各列，无须先显式求逆。

\section{拼接、分块运算与克罗内克积}

\subsection{把矩阵看作若干子块}

\keyterm{矩阵拼接}（matrix concatenation）把矩阵沿行或列方向接在一起。
若 $A\in\mathbb R^{m\times n}$、$B\in\mathbb R^{m\times p}$，则水平拼接 $[A\ B]$
为 $m\times(n+p)$ 矩阵，要求行数相同；竖直拼接要求列数相同。
增广矩阵正是水平拼接的例子。反过来，按行、列分组可得到\keyterm{分块矩阵}（block matrix）。

同型矩阵采用相同的行列分组时，加法和数乘可逐块进行。
乘法则要求左矩阵的列分组与右矩阵的行分组匹配。例如
\[
  \begin{pmatrix}A_{11}&A_{12}\\A_{21}&A_{22}\end{pmatrix}
  \begin{pmatrix}B_{11}&B_{12}\\B_{21}&B_{22}\end{pmatrix}
  =\begin{pmatrix}
    A_{11}B_{11}+A_{12}B_{21}&A_{11}B_{12}+A_{12}B_{22}\\
    A_{21}B_{11}+A_{22}B_{21}&A_{21}B_{12}+A_{22}B_{22}
  \end{pmatrix}.
\]
准确地说，若 $A_{ij}$ 为 $m_i\times n_j$ 块，$B_{jk}$ 为 $n_j\times p_k$ 块，
则结果的第 $(i,k)$ 块为 $m_i\times p_k$。
此式只是将式~\eqref{eq:matrix-product-entry} 的求和指标按组拆开；
块本身也是矩阵，乘法顺序仍不能交换。

对\keyterm{分块对角矩阵}（block diagonal matrix）
$D=\operatorname{diag}(P,Q)=\begin{pmatrix}P&0\\0&Q\end{pmatrix}$，其中 $P,Q$ 为方阵，
有 $D^r=\operatorname{diag}(P^r,Q^r)$（$r$ 为非负整数）。
若 $P,Q$ 都可逆，则 $D^{-1}=\operatorname{diag}(P^{-1},Q^{-1})$。
这些公式可直接通过分块乘法验证；一般分块矩阵的逆不能逐块求逆。

\subsection{克罗内克积}

\keyterm{克罗内克积}（Kronecker product）将一个矩阵的每个元素替换为另一个矩阵的数倍。
对 $A=(a_{ij})\in\mathbb R^{m\times n}$、$B\in\mathbb R^{p\times q}$，定义
$A\otimes B=(a_{ij}B)_{i,j}\in\mathbb R^{mp\times nq}$，右侧按块理解。例如
\[
  \begin{pmatrix}1&2\\0&1\end{pmatrix}\otimes I_2
  =\begin{pmatrix}I_2&2I_2\\0_{2\times2}&I_2\end{pmatrix}
  =\begin{pmatrix}1&0&2&0\\0&1&0&2\\0&0&1&0\\0&0&0&1\end{pmatrix}.
\]
它常用于组织由多个分量组成的线性系统。
与普通乘法不同，$A\otimes B$ 不要求 $n=p$；与逐元素乘法不同，它不要求两个矩阵同型，且通常扩大尺寸。

\section{矩阵范数与运算选择}

\subsection{用一个数衡量矩阵大小}

\keyterm{矩阵范数}（matrix norm）用非负实数度量矩阵的大小。
作为矩阵空间上的范数，它满足正定性、绝对齐次性与三角不等式：
$\|A\|=0$ 当且仅当 $A=0$，$\|\lambda A\|=|\lambda|\|A\|$，
以及 $\|A+B\|\le\|A\|+\|B\|$，其中 $A,B$ 同型。
常用的以下三种范数还满足相容乘积的次乘性 $\|AB\|\le\|A\|\|B\|$。

\keyterm{弗罗贝尼乌斯范数}（Frobenius norm）把所有元素视作一个向量，取其欧几里得长度：
\[
  \|A\|_{\mathrm F}
  =\left(\sum_{i=1}^{m}\sum_{j=1}^{n}|a_{ij}|^2\right)^{1/2}
  =\sqrt{\operatorname{tr}(A^{\mathsf T}A)}.
\]
最后一个等式适用于此处的实矩阵，因为 $A^{\mathsf T}A$ 的第 $j$ 个对角元素是第 $j$ 列的平方和。
对于同型矩阵 $A,B$，$\|A-B\|_{\mathrm F}$ 可度量逐项误差的总体大小。

另外两种常用范数是\keyterm{矩阵 $1$-范数}（matrix $1$-norm）与
\keyterm{矩阵无穷范数}（matrix infinity norm）：
\[
  \|A\|_1=\max_{1\le j\le n}\sum_{i=1}^{m}|a_{ij}|,\qquad
  \|A\|_\infty=\max_{1\le i\le m}\sum_{j=1}^{n}|a_{ij}|.
\]
前者取最大绝对列和，后者取最大绝对行和，都不是最大单个元素。
对例题~\ref{ex:matrix-operations-pair} 中的 $A$，有
$\|A\|_{\mathrm F}=\sqrt6$、$\|A\|_1=3$、$\|A\|_\infty=3$。

上述范数的次乘性也可从乘法公式核查。例如
$\sum_j|(AB)_{ij}|\le\sum_k|a_{ik}|\sum_j|b_{kj}|
\le\|B\|_\infty\sum_k|a_{ik}|$，取最大行和便得无穷范数的结论；
取列和可得 $1$-范数的结论。
对弗罗贝尼乌斯范数，对每个 $\sum_k a_{ik}b_{kj}$ 使用柯西--施瓦茨不等式后再对 $i,j$ 求和即可。

\subsection{计算前的检查顺序}

先判断目标：数据逐项处理用加减、数乘或逐元素乘除；线性变换的复合用矩阵乘法；
恢复未知量用消元，满足可逆条件时也可用逆矩阵表达。
行列式用于方阵的可逆性判别，迹提取对角和，秩描述线性无关的数量，范数度量大小与误差。

随后检查\keyterm{尺寸、顺序和存在条件}：加减要求同型，乘法要求内侧尺寸一致，
求逆要求方阵可逆，逐元素除法要求每个分母非零。
最后核对输出类型与尺寸，并尽可能代回原式验证。
这些检查比单独记住运算符号更能避免错误。

\end{document}
